\section{Introduction}
\label{sec:intro}

Hydraulic excavators spend a large part of their working cycle raising and lowering the boom
against the weight of the boom, arm and bucket. Most work on reducing this energy changes the
machine itself, through hybrid hydraulic architectures or energy recovery on the boom
\cite{casoli_hydraulic_2016,li_novel_2020}. A cheaper lever is the controller: for the same
hydraulic circuit, the way the boom is driven to its target decides how much work is spent on
overshoot, oscillation and braking.

Reinforcement learning (RL) can trade accuracy against energy directly through its reward, and
has been applied to excavator automation \cite{kurinov2020excavator,egli2022excavator}.
Two questions, however, decide whether such a controller is worth using on a real machine.
First, it must be compared with a well-tuned conventional controller on the same test
conditions, not only with itself. Second, the comparison must hold when the model has the
machine's real parameters: a policy that looks robust on a rough model may lose its advantage
once the model is realistic.

This paper studies both questions for the boom of a 30-tonne-class Doosan DX300 excavator, using
Soft Actor-Critic (SAC) \cite{haarnoja2018sac}. The contributions are:
\begin{itemize}
\item an energy--accuracy study of plain and residual SAC with domain randomisation for boom
positioning, with the energy weight chosen from a sweep and every result compared with tuned PID
baselines on 100 unseen scenarios;
\item a two-stage evaluation, first on a placeholder plant and then on a plant whose parameters
come from a Simscape model of the DX300 boom and front attachment, validated against that
model; on the placeholder plant SAC buys robustness at an energy cost, but on the real-parameter
plant a PID with an online payload estimate misses no move and uses less energy than every SAC
policy;
\item an energy-floor analysis that locates the remaining energy gap of the best controllers in
overshoot and oscillation excited by sensor noise, which sets a concrete target for future
learning-based control.
\end{itemize}

Section~\ref{sec:related} reviews related work. Section~\ref{sec:model} describes the boom model
and its two parameter sets, Section~\ref{sec:method} the learning and baseline controllers, and
Section~\ref{sec:exp} the experiments. Section~\ref{sec:results} presents the results,
Section~\ref{sec:discussion} discusses them and Section~\ref{sec:conclusion} concludes.
% "30-tonne class": trade press calls the DX300LC a 30 t excavator
% (theconstructionindex.co.uk/news/view/doosan-ups-the-spec-on-30t-excavator); confirm the
% operating weight of the quarry's variant from the Develon spec sheet.
